
%----------------------------------------------------------------------------------------
%	                            PORTADA
%----------------------------------------------------------------------------------------
% \makeportada{}{Daniel Carbajales Soto}{ 1° CFGS Automatización y Robótica Industrial}{}

%----------------------------------------------------------------------------------------
%                               TABLE OF CONTENTS
%----------------------------------------------------------------------------------------
% \maketoc
%----------------------------------------------------------------------------------------
%                                    HEADER 
%----------------------------------------------------------------------------------------
% \header{}{Daniel Carbajales Soto}{1° CFGSARI}
%----------------------------------------------------------------------------------------
%	                            TABLES
%----------------------------------------------------------------------------------------
% this is just a demo in case i need to make a table
%
%
%
% \begin{tabularx}{\linewidth}{|C{2cm}|Y|Y|C{2cm}|}
% \hline
% \makecell{Header 1} & \makecell{Header 2 \\ with linebreak} & \makecell{Header 3} & \makecell{Header 4} \\
% \hline
% Row 1 & This is a long cell text that automatically wraps in the Y column & Another long cell & Short \\
% \hline
% Row 2 & \multirow{2}{*}{Multirow text} & More text & Short \\
% \cline{1-1} \cline{3-4}
% Row 3 & & Extra text & Short \\
% \hline
% \end{tabularx}


%----------------------------------------------------------------------------------------
%	                            DOCUMENT
%----------------------------------------------------------------------------------------

\clearpage
% \fchapter{2}{Ej 5}
\phantomsection
\addcontentsline{toc}{chapter}{Ej 5}
\section*{\boldit{\textcolor{waveBlue2}{Ejercicio 5}}}Dentro de los objetivos de desarrollo sostenible, se hace mención a la capacidad de resiliencia que deben tener, entre otros, las ciudades. ¿Conocés el significado de resiliencia? ¿Te considerás una persona resiliente? ¿Qué cualidades crees que precisa una ciudad para ser resiliente?

\fsubsection{Definición de Resiliencia}
Según la ONU y el PNUD, la resiliencia es la capacidad de ecosistemas, comunidades o asentamientos para adaptarse a cambios, shocks o tensiones (como desastres climáticos o crisis económicas) sin perder su estructura esencial, resistiendo impactos y acelerando la recuperación. En el contexto de la Agenda 2030, es un prerrequisito para lograr los ODS, integrando adaptación, mitigación y gestión de riesgos para un crecimiento inclusivo y sostenible.

\fsubsection{¿Me Considero una Persona Resiliente?}
Sí, me considero resiliente en medida moderada. He enfrentado desafíos como cambios laborales durante la pandemia, adaptándome mediante aprendizaje en línea y redes de apoyo, lo que me permitió crecer. Sin embargo, reconozco áreas de mejora, como el manejo del estrés emocional, para fortalecer mi adaptabilidad.

\fsubsection{Cualidades Clave para una Ciudad Resiliente}
Una ciudad resiliente evalúa, planea y actúa ante obstáculos repentinos o graduales. Sus cualidades incluyen:
\begin{itemize}
\item \textbf{Persistencia}: Anticipa impactos actuales y futuros, invirtiendo en infraestructura duradera.
\item \textbf{Adaptabilidad}: Transforma cambios en oportunidades, como integrar energías renovables en redes urbanas.
\item \textbf{Inclusividad}: Asegura equidad en procesos de gobernanza, involucrando a comunidades vulnerables.
\item \textbf{Integración}: Combina liderazgo eficaz con inversiones en buen gobierno para resiliencia socioeconómica y ambiental.
\item \textbf{Redundancia y Flexibilidad}: Diversifica recursos (e.g., múltiples fuentes de agua) y fomenta innovación colaborativa. 
\end{itemize}
Ejemplos como Montevideo o ciudades de ONU-Habitat destacan estas cualidades para mitigar riesgos climáticos.




%%%%%%%%%%%%%%%%%%%%%%%%%%%%%%%%%%%%%%%%%%%%%%%%%%%%%%%%%%
%%%%%%%%%%%%%%%%%%%% bibliography %%%%%%%%%%%%%%%%%%%%%%%%
% \nocite{*}                                 %%%%%%%%%%%%%
% \printbibliography                         %%%%%%%%%%%%%
%%%%%%%%%%%%%%%%%%%%%%%%%%%%%%%%%%%%%%%%%%%%%%%%%%%%%%%%%%
%%%%%% Last page in case there is no bibliography %%%%%%%%
% \label{Last Page}                           %%%%%%%%%%%%%
%%%%%%%%%%%%%%%%%%%%%%%%%%%%%%%%%%%%%%%%%%%%%%%%%%%%%%%%%%

